\chapter{Containerizzazione Docker}

\section{Vantaggi della Containerizzazione}

L'utilizzo di Docker garantisce:

\begin{itemize}
    \item \textbf{Riproducibilità}: ambiente identico su qualsiasi macchina
    \item \textbf{Isolamento}: nessun conflitto con librerie di sistema
    \item \textbf{Portabilità}: deploy immediato su cloud, CI/CD, workstation locali
    \item \textbf{Versionamento}: immagini taggate per rollback facili
\end{itemize}

\begin{tipbox}[Best Practice]
La containerizzazione elimina il classico problema ``works on my machine''. Ogni esecuzione avviene in un ambiente controllato e ripetibile.
\end{tipbox}

\section{Dockerfile}

Il Dockerfile utilizza un'immagine base \texttt{python:3.12-slim} per minimizzare la dimensione:

\begin{lstlisting}[language=docker, caption=Dockerfile]
FROM python:3.12-slim

WORKDIR /app

# Install system dependencies
RUN apt-get update && apt-get install -y \
    build-essential \
    && rm -rf /var/lib/apt/lists/*

# Copy requirements and install
COPY requirements.txt .
RUN pip install --no-cache-dir -r requirements.txt

# Copy source code
COPY . .

# Set python path
ENV PYTHONPATH=/app

# Default command
CMD ["python", "-m", "src.cli"]
\end{lstlisting}

\subsection{Ottimizzazioni Docker}

\begin{description}
    \item[Multi-stage non necessario] Le dipendenze Python non richiedono compilazione pesante
    \item[Cache layer] \texttt{requirements.txt} copiato prima del codice per cache efficiente
    \item[Cleanup apt] Rimozione cache \texttt{/var/lib/apt/lists/*} per immagine più leggera
    \item[--no-cache-dir] Evita cache pip inutile nel container
\end{description}

\section{Docker Compose}

Il file \texttt{docker-compose.yml} semplifica l'esecuzione:

\begin{lstlisting}[language=docker, caption=docker-compose.yml]
services:
  etl:
    build: .
    volumes:
      - ./data:/app/data
      - ./public:/app/public
    environment:
      - PYTHONUNBUFFERED=1
    command: ["tail", "-f", "/dev/null"]
\end{lstlisting}

\subsection{Volume Mounting}

Due volumi persistenti garantiscono che:

\begin{table}[h]
\centering
\begin{tabular}{@{}lp{8cm}@{}}
\toprule
\textbf{Volume} & \textbf{Scopo} \\
\midrule
\texttt{./data:/app/data} & Input XML accessibili dal container \\
\texttt{./public:/app/public} & Output Parquet persistito su host \\
\bottomrule
\end{tabular}
\caption{Volumi Docker}
\end{table}

\section{Workflow Tipico}

\begin{lstlisting}[language=bash, caption=Comandi Docker tipici]
# Build dell'immagine
docker compose build

# Parsing di tutti i file XML
docker compose run --rm etl python -m src.cli parse \
    --input data/ --workers 8

# Query interattiva
docker compose run --rm etl python -m src.cli query \
    --table aiuti --limit 10

# Export CSV
docker compose run --rm etl python -m src.cli export \
    --table aiuti --output public/exports/aiuti.csv
\end{lstlisting}

\section{CI/CD Integration}

La natura containerizzata facilita l'integrazione in pipeline CI/CD:

\begin{figure}[h]
\centering
\begin{tikzpicture}[node distance=1.2cm, scale=0.85, transform shape]
    \node[component, fill=gray!20] (git) {Git Push};
    \node[component, right=of git] (build) {Docker Build};
    \node[component, right=of build] (test) {Run Tests};
    \node[component, right=of test] (deploy) {Deploy};
    
    \draw[arrow] (git) -- (build);
    \draw[arrow] (build) -- (test);
    \draw[arrow] (test) -- (deploy);
\end{tikzpicture}
\caption{Pipeline CI/CD tipica}
\end{figure}
